The pomeron (usually denoted $\mathbb{P}$) is a theoretical, effective object in high-energy scattering processes. It is not a fundamental particle but rather a composite structure that mediates diffraction in scattering. The properties of the pomeron include no color charge, no electric charge and  zero spin. The pomeron properties lent themselves to describing infinite gluon ladders. These are usually referred to as a chain of strong interactions between quarks and gluons, including gluon self-interactions. Specifically, gluon ladders are a pictorial and mathematical description of the propagation of gluons in high-energy scattering events. Gluon ladders refer to sequences of virtual gluon exchanges that form a 'ladder-like' structure in Feynman diagrams, where each rung of the ladder represents a gluon pair (gluon-gluon interaction) that mediates the interaction between two high-energy partons (quarks or gluons) in a scattering process. These diagrams are particularly prominent in the regime of large momentum transfer and are central to understanding the behavior of partons in deep inelastic scattering (DIS) and high-energy hadronic collisions. These infinite ladders are color-singlet, as required by the pomeron description, and can still mediate diffraction via interaction with the gluons within the ladder. Specifically, protons and neutrons can interact with photons by pomeron exchange. The pomeron then acts as the mediator enabling momentum transfer between the nucleon and the photon.




